\begin{definitionbox}{Definition (Unary Arithmetic Operation)}
\begin{definition}[Unary Arithmetic Operation]
\label{def:unary-arithmetic-operation}
Let $K$ be a domain. A \emph{unary arithmetic operation} on $K$ is a function
\[
O:K\to K.
\]
\end{definition}

\begin{remark*}[Standard quantified statement]
\[
\operatorname{UnaryOperation}_K(O)
\Longleftrightarrow
\forall x\in K,\ O(x)\in K.
\]
\end{remark*}

\begin{remark*}[Interpretation]
Unary closure is the abstract form behind successor, negation, and other
single-input transformations.
\end{remark*}

\begin{dependencies}
\begin{itemize}
  \item \hyperref[def:function]{Function}
\end{itemize}
\end{dependencies}
\end{definitionbox}

\begin{definitionbox}{Definition (Binary Arithmetic Operation)}
\begin{definition}[Binary Arithmetic Operation]
\label{def:binary-arithmetic-operation}
Let $K$ be a domain. A \emph{binary arithmetic operation} on $K$ is a function
\[
O:K\times K\to K.
\]
\end{definition}

\begin{remark*}[Standard quantified statement]
\[
\operatorname{BinaryOperation}_K(O)
\Longleftrightarrow
\forall x,y\in K,\ x\mathbin{O}y\in K.
\]
\end{remark*}

\begin{remark*}[Interpretation]
Binary closure is the abstract form behind addition and multiplication.
\end{remark*}

\begin{dependencies}
\begin{itemize}
  \item \hyperref[def:function]{Function}
  \item \hyperref[def:cartesian-product-rel]{Cartesian Product}
\end{itemize}
\end{dependencies}
\end{definitionbox}

\begin{definitionbox}{Definition (Unary Closure)}
\begin{definition}[Unary Closure]
\label{def:unary-closure}
A unary operation $O$ is \emph{closed on $K$} if applying it to an element of
$K$ produces an element of $K$.
\end{definition}

\begin{remark*}[Standard quantified statement]
\[
\forall x\in K,\ O(x)\in K.
\]
\end{remark*}

\begin{remark*}[Interpretation]
The operation does not leave the domain.
\end{remark*}

\begin{dependencies}
\begin{itemize}
  \item \hyperref[def:unary-arithmetic-operation]{Unary Arithmetic Operation}
\end{itemize}
\end{dependencies}
\end{definitionbox}

\begin{definitionbox}{Definition (Binary Closure)}
\begin{definition}[Binary Closure]
\label{def:binary-closure}
A binary operation $O$ is \emph{closed on $K$} if combining two elements of
$K$ produces an element of $K$.
\end{definition}

\begin{remark*}[Standard quantified statement]
\[
\forall x,y\in K,\ x\mathbin{O}y\in K.
\]
\end{remark*}

\begin{remark*}[Interpretation]
The operation is internal to the domain.
\end{remark*}

\begin{dependencies}
\begin{itemize}
  \item \hyperref[def:binary-arithmetic-operation]{Binary Arithmetic Operation}
\end{itemize}
\end{dependencies}
\end{definitionbox}

\begin{definitionbox}{Definition (Operation Well-Definedness)}
\begin{definition}[Operation Well-Definedness]
\label{def:operation-well-definedness}
Let $\sim$ be an equivalence relation on $A$. A binary operation $O$ respects
$\sim$ if
\[
a\sim a' \land b\sim b' \Longrightarrow
a\mathbin{O}b \sim a'\mathbin{O}b'.
\]
\end{definition}

\begin{remark*}[Standard quantified statement]
\[
\forall a,a',b,b'\in A,\,
\bigl(a\sim a'\land b\sim b'\bigr)
\Longrightarrow
a\mathbin{O}b \sim a'\mathbin{O}b'.
\]
\end{remark*}

\begin{remark*}[Interpretation]
This is the condition that lets an operation descend to equivalence classes.
\end{remark*}

\begin{dependencies}
\begin{itemize}
  \item \hyperref[def:equivalence-relation]{Equivalence Relation}
\end{itemize}
\end{dependencies}
\end{definitionbox}

\begin{lemmabox}{Lemma (Left Respect)}
\begin{lemma}[Left Respect]
\label{lem:operation-left-respect}
If a binary operation respects $\sim$, then
\[
a\sim a' \Longrightarrow a\mathbin{O}b\sim a'\mathbin{O}b.
\]
\end{lemma}
\end{lemmabox}

\begin{remark*}[Standard quantified statement]
\[
\forall a,a',b\in A,\,
a\sim a' \Rightarrow a\mathbin{O}b\sim a'\mathbin{O}b.
\]
\end{remark*}

\begin{remark*}[Interpretation]
Full respect includes respect in the left coordinate.
\end{remark*}

\begin{dependencies}
\begin{itemize}
  \item \hyperref[def:operation-well-definedness]{Operation Well-Definedness}
\end{itemize}
\end{dependencies}

\begin{lemmabox}{Lemma (Right Respect)}
\begin{lemma}[Right Respect]
\label{lem:operation-right-respect}
If a binary operation respects $\sim$, then
\[
b\sim b' \Longrightarrow a\mathbin{O}b\sim a\mathbin{O}b'.
\]
\end{lemma}
\end{lemmabox}

\begin{remark*}[Standard quantified statement]
\[
\forall a,b,b'\in A,\,
b\sim b' \Rightarrow a\mathbin{O}b\sim a\mathbin{O}b'.
\]
\end{remark*}

\begin{remark*}[Interpretation]
Full respect includes respect in the right coordinate.
\end{remark*}

\begin{dependencies}
\begin{itemize}
  \item \hyperref[def:operation-well-definedness]{Operation Well-Definedness}
\end{itemize}
\end{dependencies}

\begin{theorembox}{Theorem (Induced Operation)}
\begin{theorem}[Induced Operation]
\label{thm:induced-operation-well-defined}
The operation
\[
[a]\mathbin{\overline{O}}[b]:=[a\mathbin{O}b]
\]
on $A/{\sim}$ is well-defined if and only if $O$ respects $\sim$.
\end{theorem}
\end{theorembox}

\begin{remark*}[Standard quantified statement]
\[
\operatorname{WellDefined}(\overline{O})
\Longleftrightarrow
\forall a,a',b,b'\in A,\,
\bigl(a\sim a'\land b\sim b'\bigr)
\Rightarrow
a\mathbin{O}b\sim a'\mathbin{O}b'.
\]
\end{remark*}

\begin{remark*}[Interpretation]
This theorem is the abstract engine behind quotient constructions such as
integers and rationals.
\end{remark*}

\begin{dependencies}
\begin{itemize}
  \item \hyperref[def:quotient-set]{Quotient Set}
  \item \hyperref[def:operation-well-definedness]{Operation Well-Definedness}
\end{itemize}
\end{dependencies}
